% Post-CP23 candidate; no degree-five full-image conclusion is asserted.
\subsection{Uniform supersingular five-theta geometry and four evaluations}
\label{subsec:uniform-five-theta-geometry-candidate}
\status{Candidate; complete proofs and universal certificate submitted for independent audit.}

\begin{lemma}[A cyclic order-four presentation of every supersingular genus-two curve]
\label{lem:uniform-f0-order-four-presentation-candidate}
\textup{[Candidate.]}
Over an algebraically closed field of characteristic two, every genus-two
curve of two-rank zero is isomorphic to
\[
 C_z:\quad y^2+y=a x^5+b x^3,\qquad
 a=z^2\ne0,\quad b=z^2+z+1.
\]
Choose $\beta$ with $\beta^2+\beta=z+1$. The formula
\[
       g(x,y)=(x+1,\ y+z x^2+(z+1)x+\beta)
\]
defines an automorphism fixing the unique point $W$ at infinity,
and $g^2=\iota$ is the hyperelliptic involution.
\end{lemma}
\begin{proof}
The two-rank-zero hyperelliptic normal form has one branch point;
after placing it at infinity, Artin--Schreier reduction gives
$y^2+y=x^5+\mu x^3+\nu x$ after scaling and removing the constant.
This is also the normal-form input in
Lemma~\ref{lem:p-rank-zero-genus-two-polar-stable-candidate}.
For clarity the linear term can be removed directly. Choose a root
$t$ of
\[
                 t^{16}+\mu^4t^8+\mu^2t^2+t+\nu^4=0.
\]
After $x\mapsto x+t$, remove the fourth and second powers by the
Artin--Schreier shift $y\mapsto y+A x^2+B x$, where
$A^2=t$ and $B^2=\mu t+A$. The chosen equation for $t$ implies
$B=t^4+\mu t^2+\nu$, by injectivity of squaring, so the remaining
linear coefficient is zero. Remove the new constant by a constant
Artin--Schreier shift, possible over the algebraically closed field.
Thus the model is
$y^2+y=x^5+\mu x^3$; the same normal form is explicitly recorded in
\cite[pp.~30--31]{BoxallGrantLeprevost2001}.

Choose a nonzero root $z$ of
\[
                       (z^2+z+1)^5=\mu^5z^6.
\]
This polynomial has degree ten and constant coefficient one, so such
a root exists. If $\mu=0$, both cubic coefficients are zero and any
fifth-root scaling with $\rho^5=z^2$ gives the required model.
If $\mu\ne0$, choose $\rho^5=z^2$ and adjust it by a fifth root of
unity so that $\mu\rho^3=z^2+z+1$. The displayed fifth-power
identity permits this adjustment, since cubing permutes the fifth
roots of unity. Replacing $x$ by $\rho x$ gives $C_z$.

Put $q(x)=z x^2+(z+1)x+\beta$ and $F(x)=a x^5+b x^3$.
Directly
\[
       q(x)^2+q(x)=F(x+1)+F(x),\qquad q(x+1)+q(x)=1.
\]
These identities prove that $g$ is an automorphism and its square
adds one to $y$. Its polynomial formula extends over the unique
point at infinity, which it fixes.
\end{proof}

\begin{theorem}[A nodal polarized five-self-isogeny on every supersingular Jacobian]
\label{thm:uniform-f0-five-self-isogeny-nodal-candidate}
\textup{[Candidate.]}
Let $C$ be any genus-two curve of two-rank zero over an algebraically
closed field of characteristic two. Embed $C=\Theta\subset J(C)$
using its unique Weierstrass point $W$. There is a separable isogeny
$\phi:J(C)\to J(C)$ with
\[
                \deg\phi=25,\qquad \phi^*\Theta\equiv5\Theta,
\]
whose theta image is an integral curve in $|5\Theta|$ with exactly
twenty-four ordinary nodes and no other singularities.
\end{theorem}
\begin{proof}
Use the preceding presentation and let $g$ also denote its action
on the principal Jacobian. It preserves the polarization and has
$g^2=[-1]$, so $g^\dagger=-g$. Set $\phi=[2]+g$; then
\[
                   \phi^\dagger\phi=(2-g)(2+g)=[5].
\]
It is a separable polarized isogeny of degree twenty-five.

First $\ker\phi\cap\Theta=\{W\}$. For a finite point $P$, membership
would give $2P+gP\sim3W$. But $H^0(C,\mathcal O_C(3W))$
has basis $1,x$, so $W$ is a base point of $|3W|$. The effective
divisor $2P+gP$ does not contain $W$, a contradiction.

We next bound all fibres of $\phi|_\Theta$ using a direct reduced
divisor calculation. Write $P=(p,y_0)$, set $X=x+p$, and put
\[
 \begin{split}
 c_1&=a p^4+b p^2,\\
 c_2&=c_1+z p^2+(z+1)p+\beta,\\
 v_P(x)&=y_0+c_1X+c_2X^2,\\
 \ell_P&=a p+c_2^2+a.
 \end{split}
\]
The function $v_P$ interpolates the value and first derivative of
$y$ at $P$, and its value at $p+1$ is the $y$-coordinate of $gP$.
Using $y_0^2+y_0=F(p)$ gives the exact identity
\[
 v_P(x)^2+v_P(x)+F(x)
       =X^2(X+1)\bigl(aX^2+\ell_PX+b+\ell_P\bigr).
\]
The function $y+v_P(x)$ has its sole pole of order five at $W$.
Its zero divisor is therefore $2P+gP+E_P$, where $E_P$ is the
effective finite degree-two residual divisor whose $x$-polynomial,
with leading coefficient $a$, is
\[
 U_P(x)=a x^2+\ell_Px+
                        a p^2+\ell_Pp+b+\ell_P.
\]
This statement permits multiplicities and accidental overlap with
$P$ or $gP$: the displayed factorization subtracts their prescribed
effective divisors and leaves the residual length two.

The residual divisor cannot contain a hyperelliptic pair, since
$y+v_P(x)$ cannot vanish at both points above the same finite $x$:
their $y$-coordinates differ by one. Thus $E_P$ is nonspecial.
Equivalently, its degree-two linear system has dimension zero.
The divisor equation says
\[
                  \phi(P-W)=-[E_P-2W].
\]
Consequently equal images of finite $P,Q$ imply $E_P=E_Q$ and
$U_P=U_Q$. Writing $Q=(q,y_1)$, comparison of the linear and
constant coefficients gives
\[
 \ell_Q=\ell_P=\ell,\qquad
                 (p+q)\bigl(a(p+q)+\ell\bigr)=0.
\]
For a fixed image there are at most two possible $x$-coordinates:
$p$ and $p+\ell/a$.

There cannot be two distinct points with the same $x$-coordinate
in one fibre. They would be $P$ and $\iota P=-P$ in the Jacobian,
so $2\phi(P-W)=0$. Two-rank zero gives $J[2](k)=0$, hence
$\phi(P-W)=0$, contrary to the first paragraph. The point $W$
has no other point in its fibre by that paragraph. Every geometric
fibre of $\phi|_\Theta$ therefore contains at most two points.

The restriction is birational onto its image, since otherwise the
theta divisor would be invariant under a nontrivial kernel translation.
It is an immersion because $\phi$ is etale. Its two branches at an
identified pair have equal tangent lines only if they are a
hyperelliptic pair: the Gauss map of a principal theta curve is its
canonical degree-two map. Such a pair was just excluded. Thus its
only singularities are transverse ordinary double points.
The image is symmetric and has numerical class $5\Theta$.
Its difference from $\mathcal O_J(5\Theta)$ is a symmetric
degree-zero line bundle, hence a geometric two-torsion point;
it is trivial because $f(J)=0$. Its arithmetic genus is $26$,
so normalization genus two gives exactly twenty-four nodes.
\end{proof}

\begin{lemma}[Four evaluations and two translated pairs]
\label{lem:nonproduct-three-theta-four-values-candidate}
\textup{[Candidate.]}
Let $(A,\Theta)$ be an indecomposable principally polarized surface
over an algebraically closed field of characteristic two, with
$\Theta$ symmetric. If four distinct points $P_0,P_1,P_2,P_3$
fail to impose four independent conditions on
$H^0(A,\mathcal O_A(3\Theta))$, then
\[
                    \sum_{j\ne i}P_j\in\Theta
                           \quad(0\le i\le3).
\]
For fixed nonzero $h,k\in A$, the locus of pairs $(P,Q)$ for which
the four distinct points $P,P+h,Q,Q+k$ fail to impose four
conditions is finite. The analogous assertions hold for every
Picard twist of $\mathcal O_A(3\Theta)$ after translating coordinates.
\end{lemma}
\begin{proof}
First any three distinct points impose three independent conditions.
To separate $P_2$ from $P_0,P_1$, choose
$a\in P_0-\Theta$ outside $P_2-\Theta$ and the single point
$-P_1-P_2$. Next choose $b\in P_1-\Theta$ outside
$P_2-\Theta$ and $-a-P_2+\Theta$. The latter two intersections
are finite: equality of the second pair of theta translates would
force the excluded value $a=-P_1-P_2$. With $c=-a-b$, the product
of the three theta sections vanishes at $P_0,P_1$ and not at $P_2$.
Repeating for each excluded point proves the assertion. This also
recovers Lemma~\ref{lem:ppav-three-distinct-values} directly.

We exhibit a section vanishing at three chosen points and not at
the fourth whenever their sum is outside $\Theta$. Put
$t=P_0+P_1+P_2$ and suppose $t\notin\Theta$. In the surface
$\Theta\times\Theta$ consider the nonempty effective divisor
\[
        Z_t=\{(u,v):u+v\in\Theta+t\}.
\]
The addition map $\Theta^2\to A$ is dominant, so this is a proper
divisor. None of its curve components has constant first projection,
constant second projection, or constant sum. Indeed such a component
would respectively have $u=t$, $v=t$, or $u+v=0$: equality of the
corresponding theta translates follows from containment of integral
curves and the trivial translation stabilizer of a principal theta
divisor. Each alternative forces $t\in\Theta$.

Set $a=P_0-u$, $b=P_1-v$, and $c=-a-b$. The product of the three
theta sections has divisor
$\Theta+a+\Theta+b+\Theta+c$ and lies in $|3\Theta|$, by the
theorem of the square. It vanishes at $P_0,P_1,P_2$.
The conditions that it also vanish at $P_3$ are, respectively,
\[
 \begin{split}
 u&\in\Theta\cap(\Theta+P_0-P_3),\\
 v&\in\Theta\cap(\Theta+P_1-P_3),\\
 u+v&\in(\Theta+t)\cap(\Theta+P_0+P_1+P_3).
 \end{split}
\]
These are finite sets because the points are distinct. On any
component of $Z_t$ the three relevant maps are nonconstant, so
one can avoid all three conditions. Thus $P_3$ can be separated
from the other three. Repeat with each excluded index to prove
the necessary four-sum condition.

For the indicated two pairs put
\[
             U=P+2Q+h+k,\qquad V=2P+Q+h+k.
\]
The necessary conditions say
$U,U-h\in\Theta$ and $V,V-k\in\Theta$.
Each allowed set is finite. The homomorphism
$(P,Q)\mapsto(P+2Q,2P+Q)$ has determinant $-3$; it is a finite
etale isogeny in characteristic two. Hence its inverse image of
those finite sets is finite. Finally every Picard twist of
$3\Theta$ is a translate of that bundle, since $[3]$ is surjective.
\end{proof}

\begin{remark}[Universal certificate and remaining image gate]
\label{rem:uniform-five-theta-geometry-scope-candidate}
The complete program
\path{code/nonordinary_uniform_five_self_isogeny_identities.py}
checks nine universal characteristic-two polynomial identities,
including the automorphism, the residual quadratic, its global
coefficients and the equal-fibre factorization. It passed at
20:34:32 UTC in $0.08886713300307747$ seconds; the complete result is
\path{outputs/nonordinary_uniform_five_self_isogeny_identities.json}.
An initial coefficient-domain setup failure in its final leading-term
check is retained with the exact failed program in \path{history/};
no mathematical identity failed. The theorem proves a nodal theta
five-self-isogeny for every fixed supersingular Jacobian. It does
not yet prove the full degree-five monodromy: isolation of the
node-cover inertia and all gluing columns remains a separate gate.
\end{remark}
